% =====================================================================
% P1 Appendix C --- Reproducibility.
%
% Sources: paper/tables/README.md, paper/figures/README.md,
% experiment_log.md, code/*.py entry points.  Commands are the ones used
% for the 2026-10-21 result files; flags that differ per JSON are
% recorded in paper/tables/README.md and experiment_log.md.
% =====================================================================

\section{Reproducibility}\label{app:c}

\subsection{Environment and dependencies}\label{app:c1}

All experiments use the repository virtual environment
\texttt{.venv-c1} (Python 3.10.12) and the dependencies in
\texttt{requirements-c1.txt}. The versions used to produce the reported
files are NumPy 2.2.6, SciPy 1.15.3, pandas 2.3.3, scikit-learn 1.7.2,
Matplotlib 3.10.9, and PyTorch 2.14.1. The conformal comparisons use the
\texttt{mapie} package (1.5.0) together with the in-repository
implementations of split conformal, CQR, and the branch-union rule.

\subsection{Data generation and experiment commands}\label{app:c2}

The frozen dataset is a 128-path synthetic panel (32 seeds $\times$ four
windows, 336 hourly steps) generated by
\texttt{code/generate\_dataset.py}:
\begin{verbatim}
.venv-c1/bin/python projects/p1-c1-ope/code/generate_dataset.py \
    --n-seeds 32 --seed-start 42 --windows W1 W2 W3 W4
\end{verbatim}
The result files are produced by the following entry points, in order:
\begin{verbatim}
.venv-c1/bin/python projects/p1-c1-ope/code/coverage_eval.py
.venv-c1/bin/python projects/p1-c1-ope/code/trajectory_value_check.py
.venv-c1/bin/python projects/p1-c1-ope/code/policy_value_check.py
.venv-c1/bin/python projects/p1-c1-ope/code/observable_h_check.py --z-threshold 3
.venv-c1/bin/python projects/p1-c1-ope/code/observable_h_check.py --z-threshold 2.5
.venv-c1/bin/python projects/p1-c1-ope/code/observable_h_check.py --z-threshold 4
.venv-c1/bin/python projects/p1-c1-ope/code/observable_h_check.py --sigma-mode open
.venv-c1/bin/python projects/p1-c1-ope/code/identifiability_check.py
.venv-c1/bin/python projects/p1-c1-ope/code/stress_variant.py --gamma 4
.venv-c1/bin/python projects/p1-c1-ope/code/stress_variant.py --gamma 0
.venv-c1/bin/python projects/p1-c1-ope/code/theorem2_bounds.py
.venv-c1/bin/python projects/p1-c1-ope/code/general_pomdp_check.py
\end{verbatim}
The L0 variants (naive global, Mondrian, similarity fallback,
strategy-conditional) are selected by the flags of
\texttt{coverage\_eval.py}; the exact flag combination behind each JSON is
recorded in \texttt{experiment\_log.md} and in
\texttt{paper/tables/README.md}. Tables and figures are regenerated from
the JSON files, never hand-edited. All commands below run from the repository root:
\begin{verbatim}
.venv-c1/bin/python projects/p1-c1-ope/code/build_tables.py
.venv-c1/bin/python projects/p1-c1-ope/code/make_scis_boundary.py
.venv-c1/bin/python projects/p1-c1-ope/code/make_scis_fig1.py
.venv-c1/bin/python projects/p1-c1-ope/code/make_scis_figures.py
.venv-c1/bin/python projects/p1-c1-ope/code/latex_convert_scis.py
.venv-c1/bin/python projects/p1-c1-ope/code/latex_postprocess_scis.py
.venv-c1/bin/python projects/p1-c1-ope/code/audit_numbers.py
\end{verbatim}
The SCIS Markdown source is converted to LaTeX by \texttt{latex\_convert\_scis.py},
followed by \texttt{latex\_postprocess\_scis.py} before the two SCIS PDF builds.
\texttt{audit\_numbers.py} is read-only: it recomputes the 201 derived checks and
reports \texttt{checked 201 ok 201 bad 0}.

\subsection{Number sources}\label{app:c3}

Every number in the main text resolves to a slug in
\texttt{paper/tables/m3-derived.json}; the readable slug-to-path table is
\texttt{paper/tables/\allowbreak number-sources.md}. Table~\ref{tab:c1} lists the
headline quantities and their source paths. The per-fold strict
$\hat\kappa$ values and the finite-sample bounds are persisted separately
in \texttt{results/theorem2\_bounds.json}, written by
\texttt{code/theorem2\_bounds.py}; that file also records the 14
policy$\times$method coverage and pooled-$\hat\kappa$ match flags and the
100-replicate three-sample check behind Table~A3. The real-path
\texttt{real\_data\_replay.json} records the fold-level support behind Table~B14.

\renewcommand{\thetable}{B\arabic{table}}
\renewcommand{\theHtable}{B\arabic{table}}
\setcounter{table}{10}
\begin{table}[!ht]
\centering
\small
\caption{Headline numbers and their JSON sources. Indices in
\texttt{heldout\_regime\_coverage} follow the regime order
OPEN, EXTENDED, WEEKEND, HOLIDAY, HALT.}
\label{tab:c1}
\begin{tabular}{@{}>{\raggedright\arraybackslash}p{0.36\linewidth}>{\raggedright\arraybackslash}p{0.56\linewidth}@{}}
\toprule
Quantity & Source \\
\midrule
WEEKEND naive coverage 0.8626 &
\texttt{coverage\_gap\_naive\_global.json:}\newline
\texttt{heldout\_regime\_coverage.2} \\
WEEKEND Mondrian coverage 0.9020 &
\texttt{coverage\_gap\_exp1\_mondrian.json:}\newline
\texttt{heldout\_regime\_coverage.2} \\
L1 Wide studentized within/LORO coverage &
\texttt{trajectory\_value\_check.json:summary.}\newline
\texttt{Wide\_Rebalance\_2pct.identified\_studentized.}\newline
\texttt{\{full\_window\_within\_coverage,}\newline
\texttt{full\_window\_loro\_coverage\}} \\
L2 worst-branch marginal/union coverage &
\texttt{policy\_value\_check.json:event\_conformal.}\newline
\texttt{policies.Wide\_Rebalance\_2pct.methods.}\newline
\texttt{identified\_studentized.within\_tau2.}\newline
\texttt{branch\_union\_joint.\{worst\_branch\_marginal\_}\newline
\texttt{coverage,worst\_branch\_union\_coverage\}} \\
2A/2B/Cor.~2.2 strict ranges &
\texttt{theorem2\_bounds.json:policies.Wide\_Rebalance\_}\newline
\texttt{2pct.identified\_studentized.summary.}\newline
\texttt{\{two\_a\_strict,two\_b\_strict,cor22\_strict\}} \\
$h$ specification curve &
\texttt{observable\_h\_check.json}; variants in\newline
\texttt{trajectory\_value\_check\_obs\_h*.json} \\
Stress variant $\hat\gamma_{\mathrm{st}}$ and coverage &
\texttt{stress\_variant.json} \\
General sensor-POMDP return prediction error, regime transfer, and branch recovery &
\texttt{general\_pomdp\_check.json}; main-text Figure~5 (internal
\texttt{fig7}) provenance hashes in \texttt{fig7-general-pomdp.json} \\
\bottomrule
\end{tabular}
\end{table}

\subsection{Seeds and splits}\label{app:c4}

The dataset uses seeds 42--73, one path per seed and window. L1
calibration and evaluation use disjoint path sets of 512 paths each
(128 seeds $\times$ four windows per set), distinct from the 128-path L2
event study. L2 event studies use one
freeze-to-OPEN event per path at $\tau=2$ (main) and $\tau=81$
(transfer); the five folds are assigned by
$\text{fold}=(\text{path index}\bmod 32)\bmod 5$, so all windows of a
seed share a fold. The event-study RNG seeds are 20261029, 20261030, and
20261031; the robustness table varies them. The I5 stress variant uses
128 seeds $\times$ four windows, 384 steps per window, and
$\tau\in\{2,65,81\}$ with 512 events per horizon. The general sensor-POMDP
uses calibration seeds 3000--3127, evaluation seeds 4000--4127, and 64 target
rollouts per model and path; truth and estimation random streams are disjoint.

For Theorem~2, calibration scores are the $\tau=2$ scores of the four
folds other than the evaluation fold ($n_{\mathrm{cal}}=100$ or $104$),
and test scores are the $\tau=81$ scores of the evaluation fold
($m=28$ or $24$). The two sides are disjoint by construction because
folds are path-grouped. The strict per-fold $\hat\kappa$ is the
two-sample Kolmogorov--Smirnov distance between those two score sets;
the pooled $\hat\kappa=0.1171875$ is computed on the full calibration
and test sets and is reported only as a comparison.
